\subsubsection{General Outflow Boundary Condition}
\label{SEC:BC:Outflow:General}

This boundary condition combines subsonic and supersonic outflow boundary
conditions into a single option.
%
This boundary condition takes the velocity from the left\slash interior state,
$\vec{u}_\subL$, and evaluates the magnitude of that velocity normal to the
boundary face, $v_{n,\subL}$, as
%
\[
  v_{n,\subL} = \vec{u}_\subL\cdot\normal
\]
%
where \normal{} is the unit normal for the boundary face that points outside the
computational domain.
%
The interior static temperature $T_\subL$ and $v_{n,\subL}$ can then be used to
compute an interior Mach number normal to the outflow boundary, $M_{n,\subL}$, as
%
\[
  M_{n,\subL} = \frac{v_{n,\subL}}{\sqrt{\gamma \rgas T_\subL}}
\]
%
If $M_{n,\subL} \geq 1$, the outflow is supersonic, and the
right state is evaluated the same way as \sref{SEC:BC:Outflow:Supersonic} by
fixing it to the interior state, \thatis \[ \cvVecR = \cvVecL \]
%
Otherwise, the outflow is subsonic, and the right state is evaluated based on
the subsonic outflow boundary condition (\ref{SEC:BC:Outflow:Subsonic}), except
there is no safe guard for the outflow being supersonic since that situation is
already dealt with for this boundary condition; the safe guard on the flow
becoming reversed does still exist, but again, only provides minimal protection
against very minor violations of the assumed outflow.

%###############################################################################

